\global\long\def\numTD{Projet}%
 
\global\long\def\themeTD{Méthode stochastique pour la classification de textes littéraire en français}%
 
\global\long\def\file{toto.tex}%


\begin{center}
\begin{tabular}{|p{2cm}p{14cm}|}
\hline
{\includegraphics[width=1.8cm,viewport=0 0 337 248]{../CM/images/sorbonne.png}} & \raisebox{2ex}{\begin{Large}\textbf{Programmation de Modèles Linguistiques (I)}\end{Large}}\\

%2019-2020& \raisebox{2ex}{\begin{large}\textbf{L5SOPROG L3 Sciences du Langage}\end{large}}\\
2023-2024& \raisebox{2ex}{(L5SOPROG L3 Sciences du Langage)}\\
   & \begin{large}\textbf{\numTD}\end{large} \begin{large} \textbf{\themeTD}\end{large} \\
&\\
& Caroline Koudoro-Parfait et Gaël Lejeune, Sorbonne Université \\
\hline
\end{tabular}
\end{center}


\hrule %%%%%%%%%%%%%%%%%%%%%%%%%EN-TETE%%%%%%%%%%%%%%%%%%%%%%%%%%%%%
%\renewcommand{\contentsname}{Sommaire du TD}
%\tableofcontents
%\newpage

\noindent %
\fcolorbox{red}{lightgray}{ %
\begin{minipage}[c]{16.5cm}%
 

\section*{Objectifs}
\begin{itemize}
	\item Mettre en pratique les connaissances acquisse en Python pour l'analyse et le traitement de données textuelles.
	\item Constituer un ensemble de textes littéraires en français.
	\begin{itemize}
		\item Au moins 2 auteurs différents.
	\end{itemize}
	\item Identifier et extraire les caractéristiques statistiques du corpus.
	\item Entraîner un modèle linguistique de classification à partir d'une approche stochastique.
	\item Évaluer la précision et rappel (\textit{Precision and recall})
\end{itemize}
%
\end{minipage}} 

\section{Exercice 1 : Préparation des données et entraînement (5 points)}

\vspace{0.5cm}
 Pour ce projet, vous devez développer un classificateur de textes comme fait en TD. Pour cela, vous devez commencer par construire vos propres ressources textuelles, c'est-a-dire que vous devez chercher des oeuvres littéraires de différents auteurs en considérant les points suivants :

\begin{itemize}
	\item Format \textit{utf8}.
	\item Des textes exclusivement en français.
	\item Minimum deux auteurs, chaque auteur représentera une classe.
	\item Pensez à avoir un nombre consistant de documents pour chaque classe et équilibré entre toutes les classes.
\end{itemize}

\textbf{Exception : }

Des ressources facultatives seront mises à disposition avec une pénalisation sur la note.

\vspace{0.5cm}

\subsection{Étape 1 : Lecture de fichiers}

Faite attention avec les textes à traiter, à différence de l'anglais, le français est beaucoup plus riche en signes d'accentuation et de ponctuation. Veuillez vérifier à chaque étape que ces signes ne se perdent pas à cause d'une mauvaise manipulation.

\subsubsection{Enregistrer titres des corpus}

Créez une liste nommée \textit{input\_files}, la liste doit contenir les noms des fichiers que vous venez de télécharger.

\subsubsection{Lire les corpus}

\begin{itemize}
	\item Créez deux listes : \textit{input\_texts} et \textit{labels}.
	\item Créez une boucle avec le fonction \textit{enumerate()} en lui donnant comme argument la liste \textit{input\_files}. La fonction \textit{enumerate()} renvoie le contenu de l'objet donné commer argument, ainsi que l'index de chaque élément. Affectez ces valeurs à deux variables : \textit{f} et \textit{label} respectivement.
	\item Dans une deuxième bouclé imbriquée, parcourez le contenu des deux corpus téléchargés, utilisez la fonction \textit{open()} en lui donnant \textit{f} comme argument. Pourquoi \textit{f} ?
	\item Pour chaque itération de la boucle, la fonction \textit{open()} va renvoyer ligne par ligne le contenu du fichier indiqué en argument. Appliquez le prétraitement suivant :
	\begin{itemize}
		\item Convertissez toute la ligne en minuscule;
		\item supprimez le dernier caractère a droit (normalement l'instruction de retour à la ligne) avec la fonction \textit{rstrip()};
		\item supprimez tous les signes de ponctuation
	\end{itemize}
	 et . Affectez le résultat dans une nouvelle variable appelée \textit{line}.
	\item Avec la fonction \textit{append}, rajoutez la ligne pré-traitée à la liste \textit{input\_texts} en parallèle vous devez rajouter aussi l'étiquette associé à cette ligne (\textit{label}) à la liste \textit{labels}.
\end{itemize}

%La première boucle va juste nous montrer chaque élément dans la liste  \textit{input\_files}. On utilise la fonction \textit{enumerate()} car elle retourne les éléments de la liste et sa position. Étant donné que la liste contient les noms des fichiers et que chaque fichier va représenter une classe, on va associer chaque classe avec son index comme son identificateur numérique ou son étiquette.

%La deuxième liste va lire toutes les lignes des fichiers, convertir tout en minuscules et supprimer les ponctuation avant de les enregistrer dans l'objet \textit{input\_texts}. Notez que l'enregistrement se fait ligne par ligne, et de façon parallele, on enregistre aussi l'étiquette correspondante à la ligne dans l'objet \textit{labels.}


%\begin{python}

%# collect data into lists
%input_texts = []
%labels = []

%for label, f in enumerate(input_files):
% print(f"{f} corresponds to label {label}")

% for line in open(f):
%  line = line.rstrip().lower()
%  if line:
%   # remove punctuation
%   line = line.translate(str.maketrans('', '', #\_
%   string.punctuation))

%   input_texts.append(line)
%   labels.append(label)

%\end{python}

\subsection{Étape 3 : Division de corpus}

Nous allons réserver une partie de notre corpus pour l'évaluation. Par norme, on utilise le 70\% pour l'entraînement et le 30\% pour l'évaluation. Vous pouvez utiliser n'importe quelle technique pour effectuer cette séparation, soit directement sur python, ou bien avec des commandes \textit{bash} sur linux. Pour praticité, dans cet exemple nous allons utiliser la fonction \textit{train\_test\_split()} contenue dans la bibliothèque \textit{sklearn} que nous avons importé au début. 

Implémentez la ligne de code suivante :

\begin{python}
train_text, test_text, Ytrain, Ytest = 
train_test_split(input_texts, labels)
\end{python}
Observez que nous envoyons comme arguments les lignes du texte et ses étiquettes, comme ça nous allons garder la trace qui va nous indiquer à quelle classe appartient chaque ligne, même si elles ne sont plus dans le même ordre.

\begin{itemize}
	\item \textbf{train\_text} : contient le texte d’entraînement
	\item \textbf{test\_text} : contient le texte d'évaluation
	\item \textbf{Ytrain} : contient les étiquettes associées aux lignes du texte d'entraînement
	\item \textbf{Ytrain} : contient les étiquettes associées aux lignes du texte d'évaluation
\end{itemize}

Nous allons tester si la fonction a bien effectué la tâche demandée. Essayez le code suivant :
\begin{python}
len(Ytrain), len(Ytest)
train_text[:5]
Ytrain[:5]
\end{python}
Qu'est-ce que les commandes précédentes montrent ?

\subsection{Étape 4 : \textit{mapping}}

Comme vous sûrement savez, la machine n'est pas capable de travailler avec des données linguistiques comme les mots. Nous avons besoin de transformer les mots de notre vocabulaire en valeur numériques exploitables pour la machine.
Nous devons associer alors, chaque mot du vocabulaire à une valeur entière unique. Pour ce faire, nous avons besoin d'un objet où l'on puisse enregistrer une association de genre (${token \rightarrow id}$).
Quel objet peut-on utiliser ? Nommez l'objet \textit{word2idx}\\

Pensez à une stratégie pour parcourir tout le vocabulaire, mot par mot et à chaque nouveau mot affectez \textit{word2idx}. Il doit contenir aussi une représentation pour les mots inconnus, par exemple, les mots qui sont dans le texte d'évaluation mais pas dans le texte d'entraînement, qui nous seront donc inconnu lors de l'étape d'évaluation.

Vous pouvez vous inspirer du morceau de code suivant :

\begin{python}
idx = 1
word2idx = {'<unk>': 0}
for word in vocabulaire:
 if word not known :
  word2idx.add(word, idx)
  idx++
\end{python}

\subsubsection{Stratégie}

\begin{enumerate}
	\item Créez une boucle pour parcourir les éléments de \textit{train\_text} afin d'obtenir une ligne à la fois.
	\item Pour chaque itération faite une découpage de la ligne en mots avec la fonction \textit{split()}, affectez la liste de mot renvoyées par \textit{split()} a une variable (\textit{tokens}).
	\item Dans une deuxième boucle imbriquée, parcourez mot a mot les mots dans \textit{tokens}, validez si le mot courant n'existe déjà dans \textit{word2idx}, s'il n'existe pas, rajoutez le au dictionnaire comme \textit{key} et comme \textit{value} indiquez la valeur actuelle de \textit{idx}.
	\item Ensuite, incrémentez la valeur de \textit{idx} afin d'avoir une valeur unique associée a chaque nouveau mot.
\end{enumerate}


\begin{python}
	if token not in word2idx:
	 word2idx[token] = idx
	 idx += 1
\end{python}
%\begin{python}
%# populate word2idx
% for text in train_text:
%  tokens = text.split()
%   for token in tokens:
%    if token not in word2idx:
%     word2idx[token] = idx
%     idx += 1
%\end{python}

\begin{itemize}
	\item Affichez le contenu du dictionnaire \textit{word2idx} afin de vérifier s'il est bien constitué.
	\item Affichez aussi la taille de notre vocabulaire.
\end{itemize}

\subsection{Étape 6 : Remplacement des mots par leurs représentations numériques}

Maintenant nous allons remplacer les mots dans le vocabulaire d'entraînement par leurs valeurs numériques à l'aide de l'objet \textit{word2idx}. Vous devez créer un nouvel objet pour enregistrer les remplacements effectués.

\begin{python}
	# convert data into integer format
	train_text_int = []
	test_text_int = []
	
\end{python}
\begin{enumerate}
	\item Parcourez ligne par ligne l'objet \textit{train\_text}, utilisez une boucle.
	\item Ensuite, découpez la ligne obtenue en mots, utilisez la fonction \textit{split()}.
	\item À l'aide d'une deuxième boucle imbriquée, parcourez la liste de mots renvoyée par la fonction \textit{split()}.
	\item Dans la deuxième boucle, cherchez la représentation numérique de chaque mot de la phrase analysée.
	\item Enregistrez les représentations trouvées dans une liste temporaire (doit être réinitialisé à chaque itération) dans le même ordre d'apparition des mots.
	\item À la fin de la lecture de chaque phrase, affectez la liste temporaire avec les représentation numériques à l'objet \textit{train\_text\_int}.
	\item \textcolor{red}{L'objet \textit{train\_text\_int} peut être illustré comme une matrice dont chaque ligne est une liste représentant chaque phrase du vocabulaire}.
\end{enumerate}  

%\begin{python}
%for text in train_text:
% tokens = text.split()
% line_as_int = [word2idx[token] for token in tokens]
% train_text_int.append(line_as_int)
%\end{python}

%for text in test_text:
% tokens = text.split()
% line_as_int = [word2idx.get(token, 0) for token in tokens]
% test_text_int.append(line_as_int)

%Notez que dans notre implémentation, nous avons construit une nouvelle liste avec la même architecture que \textit{train\_text} (ligne $\times$ position), en remplaçant les mots par leurs idx.

\textbf{
Répétez la même procédure pour le corpus d'évaluation en affectant la liste \textit{test\_text\_int}}

Regardez les lignes 50, 100, 150, etc de l'objet \textit{train\_text\_int} et comparez les avec les mêmes lignes de l'objet \textit{train\_text}

\section{Exercice 2 : Normalisation de données et implémentation (7 points + 3 )}

Une fois que les caractéristiques extraites du corpus sont stockées dans nos dictionnaires, nous procédons à l'étape de normalisation.

\subsection{Étape 1 : Dictionnaire \textit{initial}}

À l'aide d'une boucle, parcourez les éléments du dictionnaire \textit{initial}.
Rappelez que dans l'index \textit{key} nous avons les mots qui commencent les phrases analysées, et dans \textit{value} nous avons leur nb d'occurrences, vous devez donc modifier cette dernière valeur en lui divisant par la somme de toutes les occurrences de tous les mots. Inspirez vous de l'exemple suivant :

\begin{python}
sum_occ = sum(initial.values())
for element in initial:
 element.value = element.value / sum_occ
\end{python}

\subsection{Étape 2 : Dictionnaires \textit{first\_order} et \textit{second\_order}}

\subsubsection{fonction \textit{list2pdict()}}
Créez la fonction \textit{list2pdict}, elle recevra comme argument une liste de mots non uniques. La fonction doit renvoyer un dictionnaire avec les mots uniques de la liste placés sur l'index \textit{key}, et en \textit{value} leur nb d'occurrences \textbf{normalisés} (réutilisez la méthode de normalisation appliquée pour le premier dictionnaire).
\begin{python}
def list2pdic(listeMots):
 dictionnaire ={}
 n = len(listeMots)
 for mot in listeMots:
  dictionnaire.key = mot
  dictionnaire.value = dictionnaire.value + 1
  
 for element in dictionnaire:
  element.value = element.value / n

\end{python}

\subsubsection{Implémentation}
Regardez le code suivante et décrivez sa fonctionnalité :
\begin{python}
for t_1, ts in first_order.items():
 # replace list with dictionary of probabilities
 first_order[t_1] = list2pdict(ts)
\end{python}
\vspace{5cm}
Réutilisez le code précédent afin de normaliser aussi le dictionnaire \textit{second\_order}.

\subsection{Fonction \textit{sample\_word()}}
Créez la fonction \textit{sample\_word()} qui devra recevoir comme argument un sous dictionnaire. Prenez le code suivant comme base et complétez la fonction pour qu'elle renvoie le mot avec une valeur cumulée supérieure à $p0$.

\begin{python}
def sample_word(d):
 #Juste une valeur de référence random
 p0 = np.random.random()
 for element in d:
  if element.value.cumulées > p0:
   return element.key #Le mot
\end{python}

\subsection{Fonction \textit{generate()}}
Créez la fonction \textit{generate()} qui devra recevoir comme argument le nombre de lignes qui seront générées. À l’intérieure de la fonction procédez de la manière suivante :

\begin{enumerate}
	\item Déterminez le premier mot du nouveau texte.
	\begin{itemize}
		\item Exécutez la fonction \textit{sample\_word()}
		\item Envoyez comme argument le dictionnaire \textit{initial}
	\end{itemize}
	 
	\item Déterminez le deuxième mot en exécutant la fonction \textit{sample\_word()}. 
	\begin{itemize}
		\item Envoyez comme argument le dictionnaire \textit{first\_order}.
		\item ATTENTION ! : Envoyez pas tout le dictionnaire, rappelez que \textit{first\_order} est un dictionnaire de dictionnaires ! Alors, envoyez QUE le dictionnaire directement associé au mot précédemment déterminé.
	\end{itemize}
	\item À l'aide d'une boucle, déterminez le reste des mots en exécutant la fonction \textit{sample\_word()}. Envoyez comme argument le dictionnaire \textit{second\_order}.
	\begin{itemize}
		\item ATTENTION ! : Envoyez pas tout le dictionnaire, rappelez que \textit{second\_order} est un dictionnaire de dictionnaires ! Alors, envoyez QUE le dictionnaire directement associé au bigramme de mots précédents.
		\item la boucle se termine quand la fonction \textit{sample\_word} détermine que le mot prochain est 'END'.
	\end{itemize}
\end{enumerate}

La fonction \textit{generate()} doit répéter (avec une boucle) la procédure décrite précédemment le nombre de fois indiqué dans l'argument afin de créer le nombre de phrases souhaité.

\section{Exercice 3 : Interface classificateur (3 points)}

\section{Exercice 4 : Évaluation (2 points)}
%\input{ressources/td2_exo3.tex}


